\documentclass[10pt,a4paper]{amsart}
\usepackage{amsmath,amssymb,amsthm,graphicx}
\usepackage{verbatim}
\newenvironment{reference}{\comment}{\endcomment}
%\newenvironment{reference}{}{}
\newenvironment{slogan}{\comment}{\endcomment}
\newenvironment{history}{\comment}{\endcomment}

% For commutative diagrams we use Xy-pic
\usepackage[all]{xy}

% We use 2cell for 2-commutative diagrams.
\xyoption{2cell}
% Theorem environments.
%
\theoremstyle{plain}
\newtheorem{theorem}[subsection]{Theorem}
\newtheorem{proposition}[subsection]{Proposition}
\newtheorem{lemma}[subsection]{Lemma}

\theoremstyle{definition}
\newtheorem{definition}[subsection]{Definition}
\newtheorem{example}[subsection]{Example}
\newtheorem{exercise}[subsection]{Exercise}
\newtheorem{situation}[subsection]{Situation}

\theoremstyle{remark}
\newtheorem{remark}[subsection]{Remark}
\newtheorem{remarks}[subsection]{Remarks}

\numberwithin{equation}{subsection}

% Macros
%
\def\lim{\mathop{\mathrm{lim}}\nolimits}
\def\colim{\mathop{\mathrm{colim}}\nolimits}
\def\Spec{\mathop{\mathrm{Spec}}}
\def\Hom{\mathop{\mathrm{Hom}}\nolimits}
\def\Ext{\mathop{\mathrm{Ext}}\nolimits}
\def\SheafHom{\mathop{\mathcal{H}\!\mathit{om}}\nolimits}
\def\SheafExt{\mathop{\mathcal{E}\!\mathit{xt}}\nolimits}
\def\Sch{\mathit{Sch}}
\def\Mor{\mathop{\mathrm{Mor}}\nolimits}
\def\Ob{\mathop{\mathrm{Ob}}\nolimits}
\def\Sh{\mathop{\mathit{Sh}}\nolimits}
\def\NL{\mathop{N\!L}\nolimits}
\def\CH{\mathop{\mathrm{CH}}\nolimits}
\def\proetale{{pro\text{-}\acute{e}tale}}
\def\etale{{\acute{e}tale}}
\def\QCoh{\mathit{QCoh}}
\def\Ker{\mathop{\mathrm{Ker}}}
\def\Im{\mathop{\mathrm{Im}}}
\def\Coker{\mathop{\mathrm{Coker}}}
\def\Coim{\mathop{\mathrm{Coim}}}

% Boxtimes
%
\DeclareMathSymbol{\boxtimes}{\mathbin}{AMSa}{"02}

%
% Macros for moduli stacks/spaces
%
\def\QCohstack{\mathcal{QC}\!\mathit{oh}}
\def\Cohstack{\mathcal{C}\!\mathit{oh}}
\def\Spacesstack{\mathcal{S}\!\mathit{paces}}
\def\Quotfunctor{\mathrm{Quot}}
\def\Hilbfunctor{\mathrm{Hilb}}
\def\Curvesstack{\mathcal{C}\!\mathit{urves}}
\def\Polarizedstack{\mathcal{P}\!\mathit{olarized}}
\def\Complexesstack{\mathcal{C}\!\mathit{omplexes}}
% \Pic is the operator that assigns to X its picard group, usage \Pic(X)
% \Picardstack_{X/B} denotes the Picard stack of X over B
% \Picardfunctor_{X/B} denotes the Picard functor of X over B
\def\Pic{\mathop{\mathrm{Pic}}\nolimits}
\def\Picardstack{\mathcal{P}\!\mathit{ic}}
\def\Picardfunctor{\mathrm{Pic}}
\def\Deformationcategory{\mathcal{D}\!\mathit{ef}}
\usepackage{cleveref}
\crefname{originalreference}{original reference}{original references}
\Crefname{originalreference}{Original reference}{Original references}
\makeatletter
\newlabel{lemma-fitting-ideal-well-defined}{{0C4B}{source}{}{sourcepointer-1}{}}
\newlabel{lemma-fitting-ideal-well-defined@cref}{{[originalreference][0][]0C4B}{[1][0][]source}{}{}{}}
\makeatother
\crefname{originalreference}{original reference}{original references}
\Crefname{originalreference}{Original reference}{Original references}
\makeatletter
\newlabel{lemma-nodal-algebraic}{{0C49}{source}{}{sourcepointer-2}{}}
\newlabel{lemma-nodal-algebraic@cref}{{[originalreference][0][]0C49}{[1][0][]source}{}{}{}}
\makeatother
\crefname{originalreference}{original reference}{original references}
\Crefname{originalreference}{Original reference}{Original references}
\makeatletter
\newlabel{lemma-reduced-quotient-regular-ring-dim-2}{{0C48}{source}{}{sourcepointer-3}{}}
\newlabel{lemma-reduced-quotient-regular-ring-dim-2@cref}{{[originalreference][0][]0C48}{[1][0][]source}{}{}{}}
\makeatother
\title{Nodal quotients of a formal power series ring characterized by the (n-1)-minors of the Jacobian matrix}
\author{The Stacks Project Authors}
\date{stacks-project:curves.tex:local-snapshot-20261012; GNU FDL 1.2 or later}
\begin{document}
\maketitle
\noindent\textit{Editorial note.} All delivered mathematical text is the author's own, copied from the pinned Stacks Project chapter file sources/stacks-curves/pinned-original-source.tex; no statement or proof step was rewritten, reordered or removed.\par\medskip
\noindent\textit{Editorial note.} The counted proof is the author's complete argument: the coordinate-and-generator change freedom from lemma-fitting-ideal-well-defined, the explicit verification of (2) and (3) in the nodal normal form, the unit-pivoting normalization loop, the block-shape estimate on the minors with the Nakayama and invertible-minor exclusions, and the final reduction to the nondegenerate quadratic term are all retained verbatim.\par\medskip
\noindent\textit{Editorial note.} The label lemma-fitting-ideal names the counted statement, and the proof environment immediately following it in the source is delivered verbatim as its proof.\par\medskip
\noindent\textit{Editorial note.} The counted statement is carried by the author's own lemma environment in the chapter source.\par\medskip
\noindent\textit{Allowed standard prerequisites (editorial).}\begin{itemize}
\item Invariance of the Fitting-ideal criterion under coordinate and generator change. Same chapter, Lemma lemma-fitting-ideal-well-defined. Lets the proof freely change coordinates and regenerize the ideal I during the normalization steps, the pivot the whole argument rests on.
\item The nodal quadratic model and reduced dimension-2 quotients. Same chapter, Lemmas lemma-nodal-algebraic and lemma-reduced-quotient-regular-ring-dim-2. Supplies side (1), the completed nodal quadric, and the structure theorem that an ideal in a two-variable power series ring defining a reduced quotient is generated by one element.
\end{itemize}\medskip
\ifdefined\hypertarget\hypertarget{sourcepointer-1}{}\fi
\noindent\textit{Source reference.} Original reference 0C4B: sources/stacks-curves/pinned-original-source.tex, source lines 4282--4282. Prerequisite pointer invoked at proof line(s) 4328; it names a standard fact of the same chapter that the counted argument uses, and every step of the counted proof built on it is retained verbatim in the delivered body.\par\medskip
\ifdefined\hypertarget\hypertarget{sourcepointer-2}{}\fi
\noindent\textit{Source reference.} Original reference 0C49: sources/stacks-curves/pinned-original-source.tex, source lines 4090--4090. Prerequisite pointer invoked at proof line(s) 4317; it names a standard fact of the same chapter that the counted argument uses, and every step of the counted proof built on it is retained verbatim in the delivered body.\par\medskip
\ifdefined\hypertarget\hypertarget{sourcepointer-3}{}\fi
\noindent\textit{Source reference.} Original reference 0C48: sources/stacks-curves/pinned-original-source.tex, source lines 4034--4034. Prerequisite pointer invoked at proof line(s) 4382; it names a standard fact of the same chapter that the counted argument uses, and every step of the counted proof built on it is retained verbatim in the delivered body.\par\medskip
\begin{lemma}
\label{lemma-fitting-ideal}
Let $k$ be a field. Let $A = k[[x_1, \ldots, x_n]]$. Let
$I = (f_1, \ldots, f_m) \subset \mathfrak m_A$ be an ideal. The following
are equivalent
\begin{enumerate}
\item $k \to A/I$ is as in Lemma \ref{lemma-nodal-algebraic},
\item $A/I$ is reduced and the
$(n - 1) \times (n - 1)$ minors of the matrix
$(\partial f_j/\partial x_i)$ generate $I + \mathfrak m_A$,
\item $\text{depth}(A/I) = 1$ and the
$(n - 1) \times (n - 1)$ minors of the matrix
$(\partial f_j/\partial x_i)$ generate $I + \mathfrak m_A$.
\end{enumerate}
\end{lemma}

\begin{proof}
By Lemma \ref{lemma-fitting-ideal-well-defined}
we may change our system of coordinates and the
choice of generators during the proof.

\medskip\noindent
If (1) holds, then we may change coordinates such that
$x_1, \ldots, x_{n - 2}$ map to zero in $A/I$ and
$A/I = k[[x_{n - 1}, x_n]]/(a x_{n - 1}^2 + b x_{n - 1}x_n + c x_n^2)$
for some nondegenerate quadric $a x_{n - 1}^2 + b x_{n - 1}x_n + c x_n^2$.
Then we can explicitly compute to show that both (2) and (3) are true.

\medskip\noindent
Assume the $(n - 1) \times (n - 1)$ minors of the matrix
$(\partial f_j/\partial x_i)$ generate $I + \mathfrak m_A$.
Suppose that for some $i$ and $j$ the partial derivative
$\partial f_j/\partial x_i$ is a unit in $A$. Then
we may use the system of parameters
$f_j, x_1, \ldots, x_{i - 1}, \hat x_i, x_{i + 1}, \ldots, x_n$
and the generators
$f_j, f_1, \ldots, f_{j - 1}, \hat f_j, f_{j + 1}, \ldots, f_m$
of $I$. Then we get a regular system of parameters $x_1, \ldots, x_n$
and generators $x_1, f_2, \ldots, f_m$ of $I$.
Next, we look for an $i \geq 2$ and $j \geq 2$ such that
$\partial f_j/\partial x_i$ is a unit in $A$. If such a pair
exists, then we can make a replacement as above and assume
that we have a regular system of parameters
$x_1, \ldots, x_n$ and generators $x_1, x_2, f_3, \ldots, f_m$ of $I$.
Continuing, in finitely many steps we reach the situation where
we have a regular system of parameters
$x_1, \ldots, x_n$ and generators
$x_1, \ldots, x_t, f_{t + 1}, \ldots, f_m$ of $I$
such that $\partial f_j/\partial x_i \in \mathfrak m_A$
for all $i, j \geq t + 1$.

\medskip\noindent
In this case the matrix of partial derivatives has the following
block shape
$$
\left(
\begin{matrix}
I_{t \times t} & * \\
0 & \mathfrak m_A
\end{matrix}
\right)
$$
Hence every $(n - 1) \times (n - 1)$-minor is in $\mathfrak m_A^{n - 1 - t}$.
Note that $I \not = \mathfrak m_A$ otherwise the ideal of minors
would contain $1$. It follows that $n - 1 - t \leq 1$ because there
is an element of $\mathfrak m_A \setminus \mathfrak m_A^2 + I$ (otherwise
$I = \mathfrak m_A$ by Nakayama). Thus $t \geq n - 2$.
We have seen that $t \not = n$ above and similarly if
$t = n - 1$, then there is an invertible $(n - 1) \times (n - 1)$-minor
which is disallowed as well. Hence $t = n - 2$.
Then $A/I$ is a quotient of $k[[x_{n - 1}, x_n]]$ and
Lemma \ref{lemma-reduced-quotient-regular-ring-dim-2}
implies in both cases (2) and (3) that $I$ is generated by
$x_1, \ldots, x_{n - 2}, f$ for some $f = f(x_{n - 1}, x_n)$.
In this case the condition on the minors exactly says that the quadratic
term in $f$ is nondegenerate, i.e., $A/I$ is as in
Lemma \ref{lemma-nodal-algebraic}.
\end{proof}

\end{document}
